\itemhead{p.67 练习 · 第1题（教材p.67）}{XJ-B1-C03-PR-P067-Y0874}
\label{XJ-B1-C03-PR-P067-Y0874}
{\footnotesize 湘教版教材p.67 / PDF第74页 / p.67 练习 / 第1题}
\field{原题}{判断下列对应关系 $f$ 是否为从 $A$ 到 $B$ 的函数：
\subq{1}$A=\{1,2,3,4\}$，$B=\{0,2,4,6,8\}$，$f$：对任意的 $x\in A$，$x\mapsto2x$；
\subq{2}$A=B=[0,+\infty)$，$f$：对任意的 $x\in A$，$x\mapsto\dfrac1x$；
\subq{3}$A=[0,+\infty)$，$B=\mathbb R$，$f$：对任意的 $x\in A$，$x\mapsto|x|$。}
\field{参考答案}{（1）是；（2）否，$0\in A$ 却没有对应值；（3）是。目标集合可以大于实际值域。}
\field{核对说明}{版本插题变化：答案册PDF19从定义域题开始，现教材先有判断对应是否为函数的一题，其余题号顺延。}
{\small \field{关联知识}{F01 函数本质、构成要素与相同函数（教材帮p.77、78 / PDF86、87）}
\field{方法与易错点}{逐一核对每个允许输入是否有且仅有一个合乎目标集合的输出。 易错：把值域直接当作目标集合；只看解析式外观；认为常函数不是函数。}
\field{答案依据}{答案册相关题组：PDF第19页（册内页18）；按当前题库原题逐项重算，不按旧版题号机械对套。}}
